\documentclass[11pt,a4paper]{article}

% Packages
\usepackage[utf8]{inputenc}
\usepackage[T1]{fontenc}
\usepackage{amsmath,amssymb,amsfonts}
\usepackage[margin=1in]{geometry}
\usepackage{enumitem}
\usepackage{hyperref}
\usepackage{fancyhdr}
\usepackage{titlesec}
\usepackage{tocloft}

% Page style
\pagestyle{fancy}
\fancyhf{}
\fancyhead[L]{\small Lecture Notes: Propositional Logic}
\fancyhead[R]{\small \thepage}
\renewcommand{\headrulewidth}{0.4pt}

% Section formatting
\titleformat{\section}{\Large\bfseries}{Section \thesection}{1em}{}
\titleformat{\subsection}{\large\bfseries}{\thesubsection}{1em}{}

% Adjust table of contents appearance
\renewcommand{\cftsecfont}{\bfseries}
\renewcommand{\cftsubsecfont}{\normalfont}

\begin{document}

% Title page
\begin{titlepage}
    \centering
    \vspace*{4cm}
    
    {\Huge\bfseries Lecture Notes:\\ Propositional Logic\par}
    \vspace{1cm}
    {\Large Discrete Mathematics\par}
    \vspace{2cm}
    {\large \textbf{Author:} Bakdaulet Sotsial\par}
    \vspace{1cm}
    
    \vfill
    {\small Astana IT University \par}
    \vspace*{2cm}
\end{titlepage}

% Table of contents
\tableofcontents
\newpage

% Section 1
\section{Propositional Logic}

\textbf{Definition (Proposition).} A \textit{proposition} (or \textit{statement}) is a declarative sentence that is either true or false, but not both simultaneously. The truth value of a proposition is its classification as true (T or 1) or false (F or 0).

\textbf{Remark.} In classical propositional logic, we adhere to the \textit{Law of the Excluded Middle}: every proposition is either true or false; and the \textit{Law of Non-Contradiction}: no proposition is both true and false.

\begin{example}
The following are propositions:
\begin{enumerate}[label=(\arabic*)]
    \item ``Paris is the capital of France.'' (True)
    \item ``$2 + 2 = 5$.'' (False)
    \item ``The number $7$ is prime.'' (True)
\end{enumerate}
\end{example}

\begin{example}
The following are \textbf{not} propositions:
\begin{enumerate}[label=(\arabic*)]
    \item ``What time is it?'' (Interrogative sentence)
    \item ``Close the door.'' (Imperative sentence)
    \item ``$x + 3 = 7$'' (Contains a free variable; truth depends on $x$)
    \item ``This statement is false.'' (Paradoxical; neither true nor false)
\end{enumerate}
\end{example}

We typically denote propositions by lowercase letters $p$, $q$, $r$, $\dots$ called \textit{propositional variables}. The area of logic dealing with propositions and their combinations is called \textit{propositional calculus}.

% Section 2
\section{Logical Operators}

Propositions can be combined using \textit{logical connectives} (or \textit{logical operators}) to form \textit{compound propositions}. The truth value of a compound proposition depends solely on the truth values of its constituent propositions and the connectives used.

\subsection{Negation (NOT)}

\textbf{Definition (Negation).} Let $p$ be a proposition. The \textit{negation} of $p$, denoted $\neg p$ (or $\sim p$, $\overline{p}$), is the proposition ``not $p$'' which is true when $p$ is false, and false when $p$ is true.

\begin{center}
\begin{tabular}{|c|c|}
\hline
$p$ & $\neg p$ \\
\hline
T & F \\
F & T \\
\hline
\end{tabular}
\end{center}

\subsection{Conjunction (AND)}

\textbf{Definition (Conjunction).} Let $p$ and $q$ be propositions. The \textit{conjunction} of $p$ and $q$, denoted $p \land q$, is the proposition ``$p$ and $q$'' which is true exactly when both $p$ and $q$ are true.

\subsection{Disjunction (OR)}

\textbf{Definition (Disjunction).} Let $p$ and $q$ be propositions. The \textit{disjunction} (inclusive or) of $p$ and $q$, denoted $p \lor q$, is the proposition ``$p$ or $q$'' which is true when at least one of $p$ or $q$ is true. It is false only when both are false.

\subsection{Implication (IF-THEN)}

\textbf{Definition (Implication).} Let $p$ and $q$ be propositions. The \textit{implication} (or \textit{conditional}) $p \to q$ is the proposition ``if $p$, then $q$'' which is false only when $p$ is true and $q$ is false. Here, $p$ is the \textit{hypothesis} (or \textit{antecedent}) and $q$ is the \textit{conclusion} (or \textit{consequent}).

\textbf{Remark.} The truth value of $p \to q$ when $p$ is false is defined as true (vacuous truth). This is a standard convention in mathematics.

\subsection{Biconditional (IF AND ONLY IF)}

\textbf{Definition (Biconditional).} Let $p$ and $q$ be propositions. The \textit{biconditional} $p \leftrightarrow q$ is the proposition ``$p$ if and only if $q$'' which is true exactly when $p$ and $q$ have the same truth value.

\textbf{Summary of Symbol Notation:}
\begin{itemize}[label=--]
    \item Negation: $\neg p$ (alternatively $\sim p$)
    \item Conjunction: $p \land q$
    \item Disjunction: $p \lor q$
    \item Implication: $p \to q$ (or $p \supset q$)
    \item Biconditional: $p \leftrightarrow q$ (or $p \equiv q$)
\end{itemize}

% Section 3
\section{Truth Tables}

A \textit{truth table} is a systematic table that displays all possible truth values of a compound proposition based on every possible combination of truth values of its atomic components. If a compound proposition involves $n$ distinct propositional variables, the truth table has $2^n$ rows.

\subsection{Truth Tables for Basic Operators}

For two variables $p$ and $q$:

\begin{center}
\begin{tabular}{|c|c|c|c|c|c|}
\hline
$p$ & $q$ & $p \land q$ & $p \lor q$ & $p \to q$ & $p \leftrightarrow q$ \\
\hline
T & T & T & T & T & T \\
T & F & F & T & F & F \\
F & T & F & T & T & F \\
F & F & F & F & T & T \\
\hline
\end{tabular}
\end{center}

\subsection{Step-by-Step Examples}

\begin{example}
Construct the truth table for $\neg p \lor q$.

First, we list all truth values for $p$ and $q$. Then compute $\neg p$, and finally $\neg p \lor q$.
\begin{center}
\begin{tabular}{|c|c|c|c|}
\hline
$p$ & $q$ & $\neg p$ & $\neg p \lor q$ \\
\hline
T & T & F & T \\
T & F & F & F \\
F & T & T & T \\
F & F & T & T \\
\hline
\end{tabular}
\end{center}
Note that this table coincides with the truth table of $p \to q$. This is an important equivalence: $p \to q \equiv \neg p \lor q$.
\end{example}

\begin{example}
Construct the truth table for $(p \lor q) \land \neg (p \land q)$. This is the \textit{exclusive or} (XOR).
\begin{center}
\begin{tabular}{|c|c|c|c|c|c|}
\hline
$p$ & $q$ & $p \lor q$ & $p \land q$ & $\neg(p \land q)$ & $(p \lor q) \land \neg(p \land q)$ \\
\hline
T & T & T & T & F & F \\
T & F & T & F & T & T \\
F & T & T & F & T & T \\
F & F & F & F & T & F \\
\hline
\end{tabular}
\end{center}
\end{example}

% Section 4
\section{Applications of Propositional Logic}

Propositional logic is a powerful tool for modeling reasoning, designing computer circuits, and formalizing specifications.

\subsection{Translating English Sentences}

To translate an English sentence into a logical expression, identify atomic propositions and the logical connectives that relate them.

\begin{example}
``If it rains today, then I will stay at home and read a book.''
\begin{itemize}
    \item Let $p$: ``It rains today.''
    \item Let $q$: ``I will stay at home.''
    \item Let $r$: ``I will read a book.''
\end{itemize}
Logical expression: $p \to (q \land r)$.
\end{example}

\begin{example}
``You can access the internet from campus only if you are a computer science major or you are not a freshman.''
\begin{itemize}
    \item Let $a$: ``You can access the internet from campus.''
    \item Let $m$: ``You are a computer science major.''
    \item Let $f$: ``You are a freshman.''
\end{itemize}
The word ``only if'' indicates the consequent: $a \to (m \lor \neg f)$. (Equivalently: if you access, then you are a CS major or not a freshman.)
\end{example}

\subsection{Real-World Applications}

\begin{itemize}
    \item \textbf{Digital Circuit Design:} Logic gates directly correspond to logical operators (AND gate $\leftrightarrow \land$, OR gate $\leftrightarrow \lor$, NOT gate $\leftrightarrow \neg$). Boolean algebra is propositional logic applied to binary signals.
    \item \textbf{Computer Programming:} Conditional statements (\texttt{if-else}) and loop conditions heavily rely on propositional logic.
    \item \textbf{Legal Reasoning and Contracts:} Contractual clauses can be formalized to check consistency and identify contradictions.
    \item \textbf{Puzzles and Games:} Many logic puzzles (e.g., knights and knaves) are solved using propositional logic.
\end{itemize}

% Section 5
\section{Propositional Equivalences}

\textbf{Definition (Tautology, Contradiction, Contingency).}
\begin{itemize}
    \item A compound proposition that is always true, regardless of the truth values of its variables, is called a \textit{tautology}.
    \item A compound proposition that is always false is a \textit{contradiction}.
    \item A compound proposition that is neither a tautology nor a contradiction is a \textit{contingency}.
\end{itemize}

\textbf{Definition (Logical Equivalence).} Two compound propositions $P$ and $Q$ are \textit{logically equivalent}, denoted $P \equiv Q$, if $P \leftrightarrow Q$ is a tautology. Equivalently, $P$ and $Q$ have identical truth tables.

\subsection{Fundamental Logical Equivalences}

Below is a list of important logical equivalences (laws). Here, $\mathbf{T}$ denotes any tautology and $\mathbf{F}$ any contradiction.

\begin{enumerate}[label=\textbf{\arabic*.}, leftmargin=*]
    \item \textbf{Identity Laws:} 
    \begin{align*}
        p \land \mathbf{T} &\equiv p, & p \lor \mathbf{F} &\equiv p
    \end{align*}
    \item \textbf{Domination Laws:}
    \begin{align*}
        p \lor \mathbf{T} &\equiv \mathbf{T}, & p \land \mathbf{F} &\equiv \mathbf{F}
    \end{align*}
    \item \textbf{Idempotent Laws:}
    \begin{align*}
        p \lor p &\equiv p, & p \land p &\equiv p
    \end{align*}
    \item \textbf{Double Negation Law:}
    \begin{align*}
        \neg(\neg p) &\equiv p
    \end{align*}
    \item \textbf{Commutative Laws:}
    \begin{align*}
        p \lor q &\equiv q \lor p, & p \land q &\equiv q \land p
    \end{align*}
    \item \textbf{Associative Laws:}
    \begin{align*}
        (p \lor q) \lor r &\equiv p \lor (q \lor r), & (p \land q) \land r &\equiv p \land (q \land r)
    \end{align*}
    \item \textbf{Distributive Laws:}
    \begin{align*}
        p \lor (q \land r) &\equiv (p \lor q) \land (p \lor r) \\
        p \land (q \lor r) &\equiv (p \land q) \lor (p \land r)
    \end{align*}
    \item \textbf{De Morgan's Laws:}
    \begin{align*}
        \neg(p \land q) &\equiv \neg p \lor \neg q \\
        \neg(p \lor q) &\equiv \neg p \land \neg q
    \end{align*}
    \item \textbf{Absorption Laws:}
    \begin{align*}
        p \lor (p \land q) &\equiv p, & p \land (p \lor q) &\equiv p
    \end{align*}
    \item \textbf{Negation Laws:}
    \begin{align*}
        p \lor \neg p &\equiv \mathbf{T}, & p \land \neg p &\equiv \mathbf{F}
    \end{align*}
\end{enumerate}

\textbf{Theorem (Implication as Disjunction).}
\begin{align*}
    p \to q \equiv \neg p \lor q
\end{align*}

\textbf{Theorem (Contrapositive).}
\begin{align*}
    p \to q \equiv \neg q \to \neg p
\end{align*}

\subsection{Worked Examples}

\begin{example}
Show that $\neg(p \lor (\neg p \land q))$ is logically equivalent to $\neg p \land \neg q$.

\textbf{Solution.} We apply a series of logical equivalences:
\begin{align*}
    \neg(p \lor (\neg p \land q)) &\equiv \neg p \land \neg(\neg p \land q) \quad \text{(De Morgan)} \\
    &\equiv \neg p \land (\neg(\neg p) \lor \neg q) \quad \text{(De Morgan)} \\
    &\equiv \neg p \land (p \lor \neg q) \quad \text{(Double Negation)} \\
    &\equiv (\neg p \land p) \lor (\neg p \land \neg q) \quad \text{(Distributive)} \\
    &\equiv \mathbf{F} \lor (\neg p \land \neg q) \quad \text{(Negation Law)} \\
    &\equiv \neg p \land \neg q \quad \text{(Identity Law)}.
\end{align*}
Thus, the equivalence holds.
\end{example}

\begin{example}
Prove that $(p \to q) \land (p \to r) \equiv p \to (q \land r)$.

\textbf{Solution.}
\begin{align*}
    (p \to q) \land (p \to r) &\equiv (\neg p \lor q) \land (\neg p \lor r) \quad \text{(Implication as disjunction)} \\
    &\equiv \neg p \lor (q \land r) \quad \text{(Distributive)} \\
    &\equiv p \to (q \land r) \quad \text{(Implication as disjunction)}.
\end{align*}
\end{example}

% Section 6
\section{Disjunctive Normal Form (DNF)}

\textbf{Definition (Literal, Minterm, DNF).}
\begin{itemize}
    \item A \textit{literal} is a propositional variable or its negation (e.g., $p$, $\neg q$).
    \item A \textit{minterm} (or \textit{elementary conjunction}) is a conjunction of literals where each variable appears at most once (e.g., $p \land \neg q \land r$).
    \item A proposition is in \textit{Disjunctive Normal Form (DNF)} if it is a disjunction of minterms.
\end{itemize}

\textbf{Definition (Full DNF).} A DNF is \textit{full} (or \textit{canonical}) with respect to a given set of variables if every minterm contains exactly one occurrence of each variable (either as the variable or its negation).

\subsection{Algorithm to Convert a Proposition to DNF}

Given a compound proposition, we can construct an equivalent DNF using the following steps:

\begin{enumerate}[label=\textbf{Step \arabic*:}]
    \item Eliminate all occurrences of $\to$ and $\leftrightarrow$ using the equivalences:
    \begin{align*}
        p \to q &\equiv \neg p \lor q \\
        p \leftrightarrow q &\equiv (p \land q) \lor (\neg p \land \neg q) \quad \text{or} \quad (\neg p \lor q) \land (p \lor \neg q)
    \end{align*}
    \item Move negations inward using \textbf{De Morgan's Laws} until negations apply only to atomic propositions:
    \begin{align*}
        \neg(p \land q) &\equiv \neg p \lor \neg q \\
        \neg(p \lor q) &\equiv \neg p \land \neg q
    \end{align*}
    \item Eliminate double negations: $\neg(\neg p) \equiv p$.
    \item Apply the \textbf{Distributive Law} of $\land$ over $\lor$ to expand the expression into a disjunction of conjunctions:
    \begin{equation*}
        p \land (q \lor r) \equiv (p \land q) \lor (p \land r).
    \end{equation*}
\end{enumerate}

\textbf{Another method (Truth Table Method):} For each row of the truth table where the proposition evaluates to T, create a minterm that is T exactly for that row (using the variable if its value is T, and its negation if F). The disjunction of all such minterms is the full DNF.

\begin{example}
Convert $p \to (q \land r)$ into DNF.

\textbf{Solution using algebraic method:}
\begin{align*}
    p \to (q \land r) &\equiv \neg p \lor (q \land r) \quad \text{(Step 1)}.
\end{align*}
This expression is already a disjunction of terms: the first term is the literal $\neg p$ (which is a minterm over $\{p\}$) and the second term is $q \land r$. To write it as a proper DNF over variables $\{p,q,r\}$, we can expand:
\begin{align*}
    \neg p \lor (q \land r) &\equiv (\neg p \land (q \lor \neg q) \land (r \lor \neg r)) \lor ((p \lor \neg p) \land q \land r) \\
    &\equiv (\neg p \land q \land r) \lor (\neg p \land q \land \neg r) \lor (\neg p \land \neg q \land r) \lor (\neg p \land \neg q \land \neg r) \\
    &\quad \lor (p \land q \land r) \lor (\neg p \land q \land r).
\end{align*}
After removing duplicates (since $A \lor A \equiv A$), the full DNF is:
\begin{align*}
    (p \land q \land r) \lor (\neg p \land q \land r) \lor (\neg p \land q \land \neg r) \lor (\neg p \land \neg q \land r) \lor (\neg p \land \neg q \land \neg r).
\end{align*}
\end{example}

\begin{example}
Find the DNF of $\neg(p \lor q) \leftrightarrow (p \land q)$.

\textbf{Solution.} Truth table method:
\begin{center}
\begin{tabular}{|c|c|c|c|c|}
\hline
$p$ & $q$ & $\neg(p \lor q)$ & $p \land q$ & $\neg(p \lor q) \leftrightarrow (p \land q)$ \\
\hline
T & T & F & T & F \\
T & F & F & F & T \\
F & T & F & F & T \\
F & F & T & F & F \\
\hline
\end{tabular}
\end{center}
The expression is true only in rows 2 ($p=T, q=F$) and 3 ($p=F, q=T$). Thus:
\begin{equation*}
    \text{DNF} \equiv (p \land \neg q) \lor (\neg p \land q).
\end{equation*}
(This is precisely the exclusive or, XOR.)
\end{example}

\textbf{Remark.} Every compound proposition has an equivalent DNF (and also an equivalent Conjunctive Normal Form, CNF). DNF is particularly useful for implementations in digital circuit design using AND and OR gates.

\end{document}